\section{Discussion and limitations}
\label{sec:limitations}

In contrast to SeqSLAM \cite{milfordSeqSLAMVisualRoutebased2012}, temporal consistency alone is insufficient for sonar: a short segment of one corridor can genuinely produce an echo sequence that closely resembles a segment elsewhere. The main remaining vulnerability therefore occurs early in a drive, when the relative pose uncertainty between two places may still span several meters and a strong aliased match remains geometrically plausible. Postponing all large corrections would reduce this risk, but would also delay the loop closures needed to constrain the map. Our ablations expose this trade-off directly: requiring 24 supporting pairs eliminated the remaining alias for the \SI{-20}{\deci\bel} sensor, but left the map of route 7 rotated. For a purely topological representation, the stricter criterion is therefore preferable, as the global rotation of the map is not that important for practical applications. For metric SLAM, however, timely anchoring of the map is also important, and a tradeoff can be made here.

A second limitation is that most experiments were conducted in simulation. The simulator models sensor directivity, geometric spreading, atmospheric absorption, and reflectivity, but does not include weak higher-order reflections. The real-world experiment demonstrates that the approach transfers to an actual acoustic sensor, requiring just the evidence curve to be recalibrated. 

The current pose graph assumes that odometry errors are predominantly random. Under a systematic heading bias, the inferred topology remains correct, but metric accuracy progressively deteriorates as expected. Introducing an explicit heading-bias state into the factor graph would provide a natural way to estimate and compensate for such systematic errors. Similarly, locations revisited only in the opposite direction are currently never associated. This limitation could be addressed by explicitly modeling how a place is expected to appear acoustically when traversed from the reverse direction, but requires complex learned world models, which is part of our future work.

Finally, although the simulated environments were deliberately designed to contain substantial self-similarity, the sequence verifier rejected nearly all incorrect hypotheses before they reached the back-end. Consequently, the GNC audit did not alter any final map in the experiments reported here. Its role is likely to become more important in environments with stronger structural aliasing, such as office floors containing repeated rooms or parking garages with highly repetitive geometry, where several individually plausible but mutually inconsistent loop-closure hypotheses may coexist. However, the target application for an algorithm like BatSLAM is in environments where optical techniques tend to fail, such as in heavy industry applications (mining, agriculture, etc). In these environments, the type of degenerative aliassing as encountered in long office corridors is typically not encountered, making this issue less pressing. In addition, all parameters were fixed using a single development trajectory, whereas the real sensor required different evidence constants. These constants primarily determine loop-closure coverage rather than map safety, but currently need to be calibrated once for each sensor configuration. Self-learning using active inference of these parameters might be an interesting avenue for further research into this architecture.

\section{Conclusion}
\label{sec:conclusion}

In this paper, we presented BatSLAM 2.0, a sonar-only SLAM system that combines a biomimetic acoustic front-end with sequence-verified place recognition and a robust pose-graph back-end. The central challenge is the strong perceptual aliasing inherent to sonar. BatSLAM 2.0 addresses this ambiguity through an explicit direction cue in the place descriptor, a sequence verifier that accepts only long, strong, unambiguous, and geometrically plausible hypotheses, a commit rule that increases the required evidence with the magnitude of the implied correction, and a back-end in which each accepted loop closure remains represented as a removable group of weak constraints.

Across all \NdesRuns{} simulated runs with calibrated or approximately calibrated odometry, the system produced consistent maps without incorrect loop closures, achieving \ATEalignedDesRange{} after alignment. When systematic odometry bias was increased beyond \SI{0.3}{\degree\per\meter}, performance degraded primarily in metric accuracy rather than in topological correctness (ie, no map collapses occured). Our analysis further showed that sensor signal-to-noise ratio and spectral-shape information are central to reliable place recognition, and that the effective capture region of a stored template is limited to approximately \CaptureLateral{} laterally and \CaptureHeading{} in heading. On a real acoustic recording, BatSLAM 2.0 likewise produced a consistent map without incorrect loop closures while operating in real time, reducing the trajectory error from \RealATEodo{}~\si{\meter} to \RealATEEven{}~\si{\meter}.

Future work will focus on substantially longer real-world trajectories using physical odometry and a bat-like binaural sensing head, explicit estimation of systematic odometry errors, and place recognition across opposite traversal directions.

\section*{Data and code availability}

The BatSLAM 2.0 source code, simulator, and scripts used to generate all experimental results, tables, and figures in this paper are available at \url{https://github.com/Cosys-Lab/BatSLAM2}.

\section*{Acknowledgment}

Claude Opus 5.5 (Anthropic) was used to assist with writing the code, generating the test harness, and drafting the text of this paper. The authors substantially revised the text and take full responsibility for the content of this paper.